\documentclass[12pt]{article}
\usepackage[T2A]{fontenc}
\usepackage[utf8]{inputenc}
\usepackage[russian]{babel}
\usepackage{amsmath,amssymb}
\usepackage{geometry}
\geometry{margin=2.5cm}
\usepackage{indentfirst}

\title{Стохастическая модель морфогенеза с топологической\\ регуляцией межклеточных контактов}
\author{А.~Ю.~Никитина, Е.~Р.~Громова}
\date{2024}

\begin{document}
\maketitle

\section*{Аннотация}
Предложена стохастическая модель роста тканей, в которой клетки обмениваются сигналами через динамически перестраиваемую сеть контактов. Показано, что топологические характеристики сети (кластеризация, распределение степеней) регулируют формирование пространственных паттернов, аналогичных наблюдаемым при эмбриональном развитии.

\vspace{0.5cm}
\noindent\textbf{Ключевые слова:} морфогенез; стохастические модели; межклеточные контакты; топология сети; бифуркации.

\vspace{0.3cm}
\noindent\textbf{Раздел журнала:} Математическое моделирование. \quad \textbf{Файл:} 1 из 1.

\section{Введение}
Морфогенез описывается системами реакционно-диффузионного типа. Мы расширяем модель, включая динамическую сеть межклеточных контактов.

\section{Модель}
Эволюция концентраций описывается стохастическим уравнением
\begin{equation}
d\mathbf{c} = \left[ D \nabla^2 \mathbf{c} + \mathbf{f}(\mathbf{c}) \right] dt + \sigma\, d\mathbf{W}_t,
\end{equation}
где $\mathbf{W}_t$ — винеровский процесс.

\section{Заключение}
Получены условия бифуркации, при которых однородное состояние теряет устойчивость.

\end{document}